% !PRINTED 444
\weberpage{444}
% running head: 444  Achtzehnter Abschnitt.  § 122.

Ferner ist, wenn $p$ nicht in $2D$ aufgeht, nach § 85
\[
D^{\frac{p-1}{2}} \equiv (D, p) \pmod{p},
\]
also:
\weq{6}{\sqrt{D}^p \equiv (D, p)\sqrt{D} \pmod{p}.}
Bezeichnet daher $\alpha$ wie oben eine ganze Zahl des Körpers $\Omega$:
\[
\alpha = \frac{x + y\sqrt{\Delta}}{2},
\]
also, wenn man mit $2Q$ multipliziert, $x$ für $Qx$ schreibt und mit $\alpha'$ die zu $\alpha$ konjugierte Zahl bezeichnet:
\[
2Q\alpha = x + y\sqrt{D}, \quad 2Q\alpha' = x - y\sqrt{D},
\]
und daraus folgt nach dem \gesperrt{Fermat}schen Satze und nach (6):
\weq{7}{\begin{aligned} \alpha^p &\equiv \alpha, & (D,p) &= +1,\\ \alpha^p &\equiv \alpha', & (D,p) &= -1, \end{aligned} \quad (\mathrm{mod}\ p).}
Daraus folgt nun, wenn $f(k, \sqrt{-m})$ irgend eine Zahl von der Form (4) ist, mit Anwendung des polynomischen Lehrsatzes:
\weq{8}{\begin{aligned} f(k,\sqrt{D})^p &\equiv f(k^p, \sqrt{D}), & (D,p) &= +1,\\ f(k,\sqrt{D})^p &\equiv f(k^p, -\sqrt{D}), & (D,p) &= -1, \end{aligned} \quad (\mathrm{mod}\ p),}
worin $k^p$ für $(k)^p$ steht, also die Bedeutung einer wirklichen Potenz hat.

3. Wenn $\mathfrak{P}$ ein in $p$ aufgehendes Primideal in $\mathfrak{K}$ ist, so ist die Norm von $\mathfrak{P}$ eine Potenz von $p$:
\weq{9}{N(\mathfrak{P}) = p^f.}

Der Exponent $f$ ist der \gesperrt{Grad} des Primideals $\mathfrak{P}$. Für jede beliebige ganze Zahl $\xi$ des Körpers $\mathfrak{K}$ ist
\weq{10}{\xi^{N(\mathfrak{P})} \equiv \xi \pmod{\mathfrak{P}},}
und wenn umgekehrt $p^f$ die niedrigste Potenz von $p$ ist, bei der die Kongruenz
\weq{11}{\xi^{p^f} \equiv \xi \pmod{\mathfrak{P}}}
durch \gesperrt{jede} Zahl $\xi$ befriedigt wird, so ist $f$ der Grad des Ideals $\mathfrak{P}$ (Bd.\,II, § 167).

Es sind nun die Grade der Primideale zu ermitteln, die in der Primzahl $p$ enthalten sind.

Da $p$ nicht in der Grundzahl des Körpers $\mathfrak{K}$ aufgeht, so ist $p$ nicht durch das Quadrat eines Primideals teilbar.
